\documentclass[12pt]{article}
\usepackage[utf8]{inputenc}
\usepackage[T1]{fontenc}
\usepackage{hyperref}
\title{Oa Language Resources for K-12 Curriculum Development}
\author{Compiled by Claude Sonnet 5 (claude-sonnet-5) for Andr\'as Kornai}
\date{}
\begin{document}
\maketitle

\section{Introduction}
Oa (ISO 639-3 tmt), also called Tasmate or Meri, is an Oceanic Austronesian language of northwest Espiritu Santo, Vanuatu. Speaker figures conflict: Wikipedia cites about 150 (1983), while Verbix gives about 300. The catalogue sources differ on its endangerment status. The search found only reference-style sources and no educational materials.

\section{Primary Educational Resources}
No primers, textbooks or school programmes were found.

\section{Linguistic and Reference Materials}
\begin{itemize}
\item Wikipedia, ``Tasmate language''. \url{https://en.wikipedia.org/wiki/Tasmate_language}
\item TLOPO database entry for Tasmate; its contents were not shown in the search. \url{https://tlopo.clld.org/languages/tasmate}
\item Endangered Languages bibliography for Tasmate. \url{https://endangeredlanguages.com/lang/3530/bibliography}
\end{itemize}

\section{Religious and Community Resources}
\begin{itemize}
\item Joshua Project profile of the Tasmate. \url{https://joshuaproject.net/people_groups/15280/NH}
\end{itemize}

\section{Digital and Audio Resources}
None confirmed.

\section{Curriculum Development Context}
Little documentation was found. The Endangered Languages bibliography is the most promising lead for locating descriptive works. Any curriculum would need new material developed with the community. Vanuatu's school language policy was not examined in this search.

\end{document}